\section{Tujuan Buku dan Capaian Pembelajaran}
Buku ajar ini dirancang untuk membimbing Anda mempelajari Pemrograman Web secara bertahap, dari konsep dasar sampai praktik pengembangan aplikasi sederhana. Pendekatan OBE memastikan setiap bab memiliki capaian yang terukur dan selaras dengan kebutuhan kompetensi program studi. Kerangka ini mengacu pada praktik penyusunan kurikulum OBE yang menekankan hasil belajar yang jelas \cite{aptikomObe}.

Melalui buku ini Anda diharapkan mampu menganalisis kebutuhan pengguna, merancang antarmuka yang aksesibel, membangun front-end dan back-end, serta melakukan evaluasi dan deployment. Porsi praktik dibuat dominan agar Anda terbiasa menerapkan konsep dalam bentuk proyek nyata. Setiap topik akan disertai contoh kode, studi kasus, dan latihan sederhana agar mudah dipahami pemula.

\begin{table}[htbp]
\centering
\begin{tabular}{|P{4cm}|P{8cm}|}
\hline
\textbf{Fokus Pembelajaran} & \textbf{Hasil yang Diharapkan} \\
\hline
Analisis kebutuhan & Mampu merumuskan kebutuhan fungsional dan non-fungsional \\
\hline
Front-end & Mampu membangun UI dengan HTML, CSS, dan JavaScript \\
\hline
Back-end & Mampu mengelola logika server dan basis data \\
\hline
Evaluasi dan deployment & Mampu menguji, mengamankan, dan merilis aplikasi \\
\hline
\end{tabular}
\caption{Fokus pembelajaran buku ajar}
\end{table}

